% q0033.tex — Drawing the line
\question{
  id        = {0033},
  title     = {Drawing the Line},
  source    = {OBECON},
  year      = {2026},
  round     = {Seletiva IEO -- Phase B},
  language  = {English},
  topic     = {Game Theory},
  subtopic  = {Mixed Strategy Nash Equilibrium / Integer Game},
  type      = {dissertative},
  statement = {%
    This question consists of two independent games that share underlying
    similarities. Read the full question before answering.

    \medskip
    \textbf{Integer Game.} Three economists simultaneously choose a positive
    integer.
    \begin{itemize}
      \item The player who chose the \emph{smallest unique} number wins.
      \item If all three choose the same number, the winner is chosen
            uniformly at random.
    \end{itemize}
    \textit{Examples}: $(1,2,3) \Rightarrow$ player with~1 wins;
    $(2,2,3) \Rightarrow$ player with~3 wins;
    $(3,3,3) \Rightarrow$ uniform random winner.

    \begin{enumerate}[(a)]
      \item (20 points)
            \begin{enumerate}[(i)]
              \item Give an example of a non-symmetric pure-strategy Nash
                    equilibrium and briefly justify why it is an equilibrium.
              \item Show that the strategy profile where each player
                    independently chooses uniformly from $\{1,2,3\}$ is
                    \emph{not} a Nash equilibrium.
            \end{enumerate}

      \item (35 points) Assume a symmetric mixed-strategy Nash equilibrium
            exists in which each player uses distribution $(p_1, p_2, p_3,
            \ldots)$ over the positive integers.
            \begin{enumerate}[(i)]
              \item Prove that the sequence $(p_n)$ is non-increasing.
              \item Prove that every $p_n > 0$.
            \end{enumerate}

      \item (20 points) \textbf{Negotiation Game.} Three firms play a
            repeated game. Each round, each firm simultaneously chooses
            \emph{Cooperate} or \emph{Defect}.
            \begin{itemize}
              \item All cooperate $\Rightarrow$ no deal, game continues.
              \item Exactly one defects $\Rightarrow$ that firm wins, game ends.
              \item Exactly two defect $\Rightarrow$ the cooperator wins, game ends.
              \item All defect $\Rightarrow$ uniform random winner, game ends.
            \end{itemize}
            \begin{enumerate}[(i)]
              \item Under any symmetric strategy, what is each firm's
                    expected payoff?
              \item If each firm defects with probability $p$, find the
                    value of $p$ that defines a symmetric mixed-strategy
                    Nash equilibrium.
            \end{enumerate}

      \item (25 points) Using the result from (c), determine the symmetric
            mixed-strategy Nash equilibrium of the Integer Game.
    \end{enumerate}
  },
  choices   = {},
  answer    = {},
  solution  = {%
    \textbf{(a.i)} $(1, 2, 2)$ is a non-symmetric pure-strategy NE. The
    player with~1 wins for sure and has no incentive to deviate; the
    players with~2 cannot improve by deviating (any other choice either
    still loses or ties in the random draw).

    \textbf{(a.ii)} Under uniform play on $\{1,2,3\}$, expected payoff is
    $1/3$. Deviating to always play~1: wins whenever neither opponent plays
    1 (probability $4/9$) plus $1/3$ chance when all play~1 ($1/9$).
    Payoff $= 4/9 + 1/27 = 13/27 > 1/3$. Profitable deviation $\Rightarrow$
    not a NE.

    \textbf{(b.i)} In a symmetric NE the winning probability is maximised
    by placing mass on smaller numbers. If $p_m > p_n$ for some $m > n$,
    swapping those probabilities increases the chance of holding the
    smallest unique number without changing payoffs in tie scenarios.
    Hence $(p_n)$ must be non-increasing.

    \textbf{(b.ii)} Suppose there is a smallest $n$ with $p_n = 0$ (and
    $n > 1$, since $p_1 = 0$ gives a trivially degenerate equilibrium).
    Redistributing half of $p_{n-1}$ to $p_n$ strictly increases expected
    payoff (new winning cases open up without losing existing ones),
    contradicting equilibrium. Hence $p_n > 0$ for all $n$.

    \textbf{(c.i)} By symmetry each firm wins with probability $1/3$,
    so expected payoff is $\boxed{1/3}$.

    \textbf{(c.ii)} Swapping ``Cooperate'' and ``Defect'' leaves the game
    structurally identical, so in a symmetric NE $p = 1-p$, giving
    $p = 1/2$.

    \textbf{(d)} The Integer Game and the Negotiation Game are structurally
    equivalent: choosing integer~$n$ corresponds to planning to defect in
    round~$n$ if no one defected earlier. Under the $p = 1/2$ equilibrium
    of the Negotiation Game, the probability of first defecting in round~$n$
    is $(1/2)^n$. Hence the unique symmetric mixed-strategy NE of the
    Integer Game is:
    \[
      p_n = \left(\tfrac{1}{2}\right)^n, \quad \forall\, n \in \mathbb{N}.
    \]
  },
}
